\documentclass[fleqn,12pt]{wlscirep}

\usepackage[utf8]{inputenc}
\usepackage[version=4]{mhchem}
\usepackage[flushleft]{threeparttable}
\usepackage{booktabs}       % professional-quality tables
\usepackage{makecell}
\usepackage{siunitx}
\usepackage{setspace}
\usepackage{array}
\usepackage{silence}
\usepackage{bm}
\usepackage{wrapfig}
\usepackage{placeins}
\usepackage{fix-cm}


\WarningFilter{microtype}{Unable to apply patch}
\newcommand{\PreserveBackslash}[1]{\let\temp=\\#1\let\\=\temp}
\newcolumntype{C}[1]{>{\PreserveBackslash\centering}p{#1}}
\newcolumntype{R}[1]{>{\PreserveBackslash\raggedleft}p{#1}}
\newcolumntype{L}[1]{>{\PreserveBackslash\raggedright}p{#1}}
\usepackage[T1]{fontenc}
\renewcommand{\thefootnote}{\fnsymbol{footnote}}

\newcommand{\WH}[1]{{\color{red}[WH: #1]}}
\newcommand{\ZD}[1]{{\color{blue}[ZD: #1]}}
\newcommand{\AP}[1]{{\color{green}[AP: #1]}}
\newcommand{\recheck}[1]{{\color{black}#1}}
% \newcommand{\recheck}[1]{{\color{red}#1}}


\title{
Smooth Equivariant Zone-bridging Model xxx
}

\author[1,2,3]{xxx}

% \author[15,16,\thanks{\href{mailto:wang_han@iapcm.ac.cn}{wang\_han@iapcm.ac.cn}}]{Han Wang}


\affil[1]{xxx}


\begin{abstract}
test text
\end{abstract}
\begin{document}

\flushbottom
\maketitle
% * <john.hammersley@gmail.com> 2015-02-09T12:07:31.197Z:
%
%  Click the title above to edit the author information and abstract
%
\thispagestyle{empty}

% \noindent Please note: Abbreviations should be introduced at the first mention in the main text – no abbreviations lists. Suggested structure of main text (not enforced) is provided below.

\section{Introduction}
 
% ---------- ¶1: motivation ----------
% TODO: ~3 sentences on DFT-grounded universal PES and MLIPs as
% near-linear-cost surrogates. Compress aggressively; this is a
% follow-up paper and readers know the background.
 
% ---------- ¶2: from DPA3 to DPA4 ----------
% TODO: 2--3 sentences sketching where DPA3 [cite] left off.
%   - DPA3 established an invariant multi-layer GNN with smooth /
%     conservative guarantees and competitive results across
%     materials and molecular benchmarks.
%   - One sentence framing what's still missing --> opens ¶3.
 
% ---------- ¶3: the three tensions DPA4 attacks ----------
% TODO: state the three problems as a paragraph or bullet list.
%   (a) Accuracy ceiling. SO(3)-equivariant models (MACE, NequIP,
%       EquiformerV2) reach high accuracy at high cost; invariant
%       models like DPA3 are fast but lag on small-molecule AIMD
%       tasks (rMD17, 3BPA), as DPA3 itself acknowledges in its
%       Discussion section.
%   (b) Compute bottleneck. DPA3 saturates the memory of an A800 80 GB at moderate model size, and the existing implementation cannot benefit from torch.compile due to dynamic shapes / in-place
%       updates / Python control flow on the hot path. There is no
%       clear way to deliver matched accuracy with less memory and
%       higher throughput within the current MLIP design space.
%   (c) Physical robustness at short range. Existing MLIPs handle
%       short-range repulsion either by data-driven extrapolation
%       (unreliable) or by post-hoc ZBL wrappers, which break the
%       smoothness or conservativeness of the underlying NN model.
 
% ---------- ¶4: contributions ----------
% TODO: state three contributions matching ¶3.
%   1. SO(2)-equivariant local-frame architecture with attention.
%      In each edge's local frame, the residual rotation symmetry is
%      SO(2), avoiding the cubic cost of SO(3) Clebsch--Gordan tensor
%      products. Result: accuracy comparable to SO(3)-equivariant
%      models at a cost comparable to invariant models. Augmented with
%      attention to capture longer-range many-body context.
%   2. End-to-end efficiency. Algorithmic gains (cheap SO(2) ops) plus
%      full torch.compile compatibility yield 5--10x training and
%      inference speedup over DPA3 at matched accuracy, with
%      substantially lower memory --- usable on commodity GPUs.
%   3. Native ZBL coupling via a C^2-smooth switch, applied per pair
%      before any energy summation, preserving conservativeness and
%      translational/rotational invariance by construction.
 
% Headline numerical claim (one sentence):
% "DPA4 [TODO: top rank on Matbench Discovery, SOTA on SPICE-MACE-OFF
% and TorsionNet-500] while training 5--10x faster than DPA3 at
% matched accuracy."
 
% ---------- ¶5: organization ----------
% TODO: one short paragraph mapping the rest of the paper.
 
% =============================================================
\section{Results}
\label{sec:results}
 
% Figure ordering: the user-suggested layout puts the efficiency Pareto
% plot as the headline figure (Fig. 1), since it encapsulates the
% strongest message in one image. Architecture diagram follows as
% Fig. 2. 
 
\subsection{The DPA4 architecture}
\label{sec:arch_overview}
 
% TODO: Fig. 2 --- main architecture figure (4 panels):
%   (a) overall block diagram
%   (b) SO(2)-equivariant local-frame edge update
%   (c) attention block
%   (d) ZBL coupling branch
%
% TODO: write 3 short paragraphs.
%
% ¶1 --- Local-frame SO(2) equivariance.
%   Why SO(2) instead of SO(3): for each edge (i, j), set the local
%   z-axis along the bond direction; the residual rotation symmetry
%   becomes the abelian group SO(2). Equivariant operations under
%   SO(2) are diagonal in the angular-order basis (e^{i m theta})
%   and avoid the cubic cost of SO(3) Clebsch--Gordan tensor products.
%   Headline: accuracy close to SO(3) models, cost close to invariant
%   models.
%   One-sentence differentiation from MACE / EquiformerV2 / eSEN /
%   EScAIP --- emphasize what makes DPA4 distinct (e.g. specific
%   choice of frame canonicalization / order of attention / fully
%   compile-friendly).
%
% ¶2 --- Edge-conditioned attention.
%   Why attention: pure local message passing struggles to capture
%   long-range many-body correlations within a single layer.
%   Motivating example: cite or describe behavior on a long-range +
%   many-body system (e.g. zeolite frameworks or OC20 catalysis
%   surfaces) where attention quantitatively closes the gap.
%   Where attention is applied: edge-level / vertex-level / cross-layer
%   --- one sentence here, full math deferred to Methods.
%   Smoothness: attention weights are masked by a smooth cutoff so the
%   resulting energy remains C^2 in atomic positions.
%
% ¶3 --- Native ZBL branch.
%   Existing approach (DPA3, MACE-MP-0): ZBL is wrapped externally,
%   adding a piecewise term that can break C^2 smoothness or
%   conservativeness of the combined potential.
%   DPA4 design: ZBL is a built-in pair branch combined with the NN
%   branch via a C^2-smooth switch, before total-energy summation.
%   Forces and virials therefore remain conservative by autograd
%   construction. Full equations deferred to Methods.
 
\subsection{Matbench Discovery}
\label{sec:matbench}
 
% TODO: Tab. 1 --- Matbench Discovery leaderboard.
%   - Reuse DPA3 Table 1 layout; place DPA4 in the top row.
%   - Required metrics: CPS, Acc, F1, DAF, Prec, MAE, R2,
%     kappa_SRME, RMSD.
%   - Compliance: trained on MPtrj only, evaluated on WBM.
%   - Headline: DPA4 achieves the highest CPS among compliant models.
 
% TODO: 1 short paragraph reporting multi-seed stability of CPS to
% pre-empt the "leaderboard gaming" reviewer concern.

\subsection{Small-molecule benchmarks}
\label{sec:small_mol}
 
% TODO: SPICE-MACE-OFF (Fig. 3 left or top panel).
%   - Same train/val/test split as Ref. [Kovacs2023] / DPA3.
%   - Three model sizes: DPA4-{S, M, L} (or whatever sizes mirror
%     DPA3-L{6,12,24}).
%   - Baselines: MACE(M/L), eSEN, EScAIP.
%   - Reuse DPA3 Fig. 2(a) layout (per-subset bar chart of energy MAE).
%   - Headline: DPA4-L surpasses eSEN (previous SOTA) by [X%] on
%     energy LWAMAE.
 
% TODO: TorsionNet-500 (Fig. 3 right or bottom panel; zero-shot from
% the SPICE-trained DPA4).
%   - Metrics: MAE, RMSE, MAEB, NABH_h.
%   - Target: NABH_h = 0 for the larger DPA4 sizes; MAEB < [TODO]
%     kcal/mol.
 
% TODO (optional, P1): rMD17 / 3BPA / acetylacetone.
%   - Direct response to DPA3's stated limitation on small-molecule
%     AIMD tasks.
%   - One small table comparing DPA4 against MACE on these AIMD
%     datasets.
%   - Drop to Supplementary if results are not strongly favorable.

\subsection{Training and inference efficiency}
\label{sec:efficiency}
 
% This subsection carries the strongest single message. Lead with the
% Pareto figure.
 
% TODO: Fig. 1 --- the headline Pareto plot.
%   - x-axis: cumulative GPU-hours (or wall-clock training time), log scale
%   - y-axis: WBM energy MAE (Matbench Discovery test), log scale
%   - Plot every Matbench-Discovery-compliant model for which training
%     trajectories are available (CHGNet, M3GNet, MACE-MP-0, SevenNet,
%     GRACE, AlphaNet, MatRIS, eSEN-30M-MP, ORB, eqV2, DPA3-L24, ...).
%   - Highlight DPA4 as a new Pareto frontier strictly below and to
%     the left of all baselines.
%   - This figure can also serve as the visual abstract.
 
% TODO: Tab. 2 --- training & inference efficiency details.
%   Rows: DPA3-L24, DPA4-L24, DPA4-L24 (no compile), MACE(L), eSEN.
%   Columns:
%     - # parameters
%     - Training throughput (samples / s)
%     - GPU-hours to reach target WBM MAE = [TODO: X] meV/atom
%     - ASE single-point latency at 80 / 200 / 500 / 1000 / 5000 atoms (ms)
%     - Peak GPU memory (GB) at fixed batch size
 
% TODO: write 3 short paragraphs.
 
% ¶1 --- Algorithmic efficiency.
%   - SO(2) ops are diagonal in the angular-order basis: O(d * m_max)
%     per edge, vs SO(3) Clebsch--Gordan products that scale as O(L^6)
%     in the maximum spherical-harmonic order.
%   - Report FLOPs ratio of DPA4's edge update vs a comparable
%     MACE-style SO(3) model.
 
% ¶2 --- Implementation efficiency (torch.compile).
%   - Why traditional MLIPs (incl. DPA3) cannot be torch.compile'd
%     cleanly: dynamic neighbor counts, in-place graph buffers, Python
%     control flow on the hot path, custom CUDA kernels with no
%     dispatcher coverage.
%   - DPA4's redesign: padded-neighbor representation, functional
%     (non-in-place) updates, use of torch.fx / make_fx-style tracing
%     where appropriate, plus full dispatcher coverage on hot ops.
%   - Report compile-time overhead and steady-state speedup from
%     enabling torch.compile alone (i.e. with/without compile on
%     identical model + hardware).
 
% ¶3 --- End-to-end claim and hardware portability.
%   - Headline: at matched WBM MAE, DPA4 reaches the target in
%     [TODO: 5--10x] fewer GPU-hours and uses [TODO: X%] of DPA3's
%     peak memory.
%   - Inference at large system size (5000 atoms): [TODO: Yx] faster
%     than DPA3-L24, [TODO: Zx] faster than MACE(L).
%   - Portability: report training throughput on V100 and RTX 4090
%     (in addition to A800) to demonstrate that DPA4 is usable on
%     commodity GPUs, not just data-center accelerators.
 
 
\subsection{Native ZBL coupling and extreme conditions}
\label{sec:zbl}
 
% TODO: Fig. 4(a) --- short-range energy curve.
%   - For one or two atom pairs, plot the combined pair energy vs r,
%     showing the smooth crossover from NN to ZBL.
%   - Annotate r_in and r_out, and verify C^2 continuity visually.
 
% TODO: Fig. 4(b) --- cascade case study.
%   - Pick the simplest example with a clear DPA4-ZBL vs DPA4-noZBL gap.
%   - Suggested: PKA cascade in W (well-known irradiation benchmark).
%   - Plot total energy drift / atomic-displacement stability vs
%     simulation time; show DPA4-noZBL diverging while DPA4-ZBL stays
%     stable.
%   - Optional: overlay DPA3-with-external-ZBL or MACE-MP-0-with-ZBL
%     to show that wrapper-style ZBL breaks conservativeness on a
%     measurable timescale.
 
% TODO: write 1--2 paragraphs.
%
% ¶1 --- Why wrapper-style ZBL is not enough.
%   Briefly recap how DPA3-style or MACE-MP-0-style ZBL is applied
%   externally and why this can break C^2 smoothness or
%   conservativeness.
%
% ¶2 --- DPA4's native coupling and impact.
%   Smooth switch is applied per pair before energy summation; forces
%   and virials remain conservative by construction. Use the case
%   study to illustrate practical impact (high-energy collisions,
%   compression/melting, rare close-distance configurations on
%   reaction paths).
 
\subsection{Architecture ablations (optional)}
\label{sec:ablation}
 
% Include only if time permits before submission. Otherwise drop.
 
% TODO: Tab. 3 --- ablation study.
%   Suggested rows (refine based on actual DPA4 design):
%     - Full DPA4
%     - DPA4 with l = 0 only (kill SO(2) features, keep attention)
%     - DPA4 without attention
%     - DPA4 without ZBL
%   Columns:
%     - SPICE-MACE-OFF energy LWAMAE
%     - MPtrj/WBM energy MAE
%     - Training throughput (samples / s)
%     - Inference latency at [TODO: X] atoms
%
% TODO: 1 paragraph interpreting the ablation results.
%   - SO(2) features and attention should each be measurably important
%     on accuracy benchmarks.
%   - ZBL has minor effect on standard benchmarks; its value shows up
%     in the extreme-condition case studies of Section 2.5.
 
% =============================================================
\section{Discussion}
\label{sec:discussion}
 
% TODO: structure as 3 paragraphs.
 
% ¶1 --- summarize the three contributions and how they interlock.
%   - SO(2) equivariance: accuracy at moderate cost.
%   - Attention: longer-range correlations within each MP step.
%   - Native ZBL: extends domain of validity to short interatomic
%     distances without breaking conservativeness.
 
% ¶2 --- limitations.
%   - DPA4 is presented here as a single-task / per-dataset model; we
%     have not yet integrated DPA3's multi-task dataset-encoding
%     scheme, nor trained a DPA4-based foundation potential. We
%     anticipate doing this in follow-up work.
%   - Charged species, open-shell magnetism, and relativistic heavy
%     elements remain out of scope.
%   - SO(2) is in principle a strict subset of SO(3) equivariance; the
%     gap appears empirically negligible but may resurface in
%     symmetry-sensitive tasks (e.g. accurate phonons in highly
%     symmetric crystals).
 
% ¶3 --- outlook.
%   - DPA4 as the base architecture for a future foundation potential
%     trained on OpenLAM-v1 / OMol25 in a multi-task setup.
%   - Downstream applications enabled by DPA4's efficiency: long-time
%     MD, irradiation-damage simulation, high-throughput catalysis
%     screening on commodity GPUs.
 
% =============================================================
\section{Methods}
\label{sec:methods}
 
\subsection{Datasets}
 
% TODO: dataset descriptions.
%   - For datasets reused from DPA3 (SPICE-MACE-OFF, TorsionNet-500,
%     MPtrj+WBM), write a one-sentence summary and cite DPA3's
%     Methods section.
%   - For any newly added dataset, give a full description.
 
\subsection{DPA4 model architecture}
 
% Describe DPA4 standalone --- do not frame it as a modification of
% DPA3's LiGS construction.
 
% TODO: notation paragraph.
%   - System with N atoms, atomic numbers Z, positions R.
%   - Energy decomposition E = sum_i E_i; forces and virials by autograd.
%   - State the DPA4 atomic energy combining the NN branch and the
%     native ZBL branch:
%       E_i^total = E_i^{NN} + sum_{j in N(i)} w_ZBL(r_ij) * E_ij^ZBL
%     where E_i^{NN} comes from the SO(2)-equivariant + attention
%     network and the second term is the smooth-switched ZBL
%     contribution.
 
\subsubsection{Local frames and SO(2)-equivariant features}
 
% TODO: define the local frame F_ij = (e1, e2, e3) for each directed
% edge (i, j).
%   - e3 = r_ij / |r_ij|.
%   - Describe the canonicalization of e1, e2 orthogonal to e3.
%     **This is your core know-how --- write the actual scheme.**
 
% TODO: define the angular-order grading on edge / vertex features.
%   - Channels labeled by m in {-m_max, ..., m_max}.
%   - Transformation rule under rotation by theta about e3:
%     phi^{(m)} -> e^{i m theta} phi^{(m)}.
 
% TODO: write the message and update equations.
%   - Edge-to-vertex message respecting angular-order grading.
%   - Frame-aligned aggregation: rotate per-edge frames into a per-vertex
%     frame F_i before summation, so global SO(3) equivariance is
%     preserved.
%   - Vertex update with residual connections and trainable step sizes.
 
\subsubsection{Attention}
 
% TODO: write the attention formula.
%   - Queries q_i, keys k_j, values v_j computed from the
%     SO(2)-invariant subspace (m = 0 channels) so attention weights
%     are rotationally invariant.
%   - softmax over neighbors j of i.
%   - Smooth cutoff applied to attention weights to preserve C^2
%     smoothness of the resulting energy.
%   - Specify exactly where in the layer attention enters:
%     edge-level / vertex-level / cross-layer (fill in based on the
%     actual DPA4 design).
 
\subsubsection{Native ZBL coupling}
 
% TODO: ZBL pair energy.
%   - E_ij^ZBL = (Z_i Z_j e^2 / 4 pi eps_0 r_ij) * Phi(r_ij / a_ZBL)
%   - Phi: universal screening function (cite Ziegler-Biersack-Littmark).
%   - a_ZBL: standard ZBL screening length.
 
% TODO: smooth switch w_ZBL(r).
%   - w_ZBL(r) = 1 for r <= r_in, 0 for r >= r_out, smooth interpolation
%     in between.
%   - State the chosen functional form (tanh / cosine / polynomial)
%     and verify C^2 continuity at r_in and r_out.
%   - Default values: [TODO: r_in = ?, r_out = ?].
 
% TODO: state that the ZBL term enters the total energy *before*
% differentiation, so forces and virials remain conservative by
% autograd construction. Contrast briefly with the wrapper-style
% application in DPA3 / MACE-MP-0.
 
\subsubsection{Symmetry guarantees}
 
% TODO: brief proof in the style of DPA3 Section 4.2.
%   - Translational invariance: trivial (relative positions only).
%   - Permutation equivariance: standard argument from sum-pooling and
%     identical message functions across same-species atoms.
%   - Rotational invariance: argue from the SO(3)-equivariance of the
%     final vertex feature, with the fitting MLP applied to its
%     SO(2)-invariant (m = 0) subspace.
%   - Conservativeness: from autograd-based force/virial definition
%     applied to the combined NN + smooth-ZBL energy.
 
\subsection{Training}
 
% TODO: hyperparameter table per benchmark (analogous to DPA3 Table S-6).
%   - Optimizer, learning-rate schedule, batch size, # GPUs, # steps.
%   - Loss function: MAE.
%   - Energy / force / virial loss prefactors.
%   - Activation function.
 
\subsection{Inference and \texttt{torch.compile}}
 
% TODO: describe the implementation work that enables full torch.compile.
%   - Padded-neighbor representation to remove dynamic shapes.
%   - Functional (non-in-place) updates of graph buffers.
%   - Use of torch.fx / make_fx-style tracing for the hot path.
%   - CUDA graph compatibility.
%   - Compile-time overhead and steady-state speedup numbers.
 
% =============================================================
\section{Acknowledgments}
 
% TODO: copy the acknowledgments block from DPA3 and add any new
% funding sources, collaborators, or computing centers used for DPA4.
 
% =============================================================
\section{Data availability}
 
% TODO: fill in URLs.
%   - Datasets used: cite original sources if reused.
%   - DPA4 source code: DeePMD-kit repository, version >= [TODO: X.Y.Z].
%   - Trained DPA4 checkpoints: [TODO: URL].
\end{document}